\documentclass[english]{article}
\usepackage[T1]{fontenc}
\usepackage[latin9]{inputenc}
\usepackage{geometry}
\geometry{verbose,tmargin=3.5cm,bmargin=4cm,lmargin=3.8cm,rmargin=3.8cm}

\makeatletter
\usepackage{url}
\usepackage{lipsum}  
\usepackage{graphicx}  

\makeatother

\usepackage{babel}
\begin{document}

\title{
    \includegraphics[width=8cm, keepaspectratio]{tudelftlogo.png}\\
    \vspace*{2cm}
    \textbf{
        Optimizing (Small) Transformer Models' Inference at the Pre-Training Stage: GQA and Dimensionalities of Keys and Values  \\
        {\large Balancing Inference Speed and Quality: Seeking Optimal Trade-offs}
    }\\
    \vspace*{1cm}
}

\author{
    \textbf{
    Khalit Gulamov$^1$}\\
    \hfill \break
    \textbf{Supervisor(s): Maliheh Izadi$^1$, Aral de Moor$^1$ }\\
    \break
%    \affiliations
    {\large 
        \hfill \break
        $^1$EEMCS, Delft University of Technology, The Netherlands
    }\\
}

\date{}

\maketitle
\thispagestyle{empty}

\let\clearpagebackup\clearpage
\renewcommand{\clearpage}{ }

\onecolumn

\vspace*{1.5cm}
\begin{center}
    A Thesis Submitted to EEMCS Faculty Delft University of Technology,\\
    In Partial Fulfilment of the Requirements\\
    For the Bachelor of Computer Science and Engineering\\
    \today
\end{center}

\vspace*{2cm}

\noindent
{\small
Name of the student:  Khalit Gulamov\\
Final project course: CSE3000 Research Project\\
Thesis committee: Maliheh Izadi, Aral de Moor, \textless Examiner\textgreater\\
}
\vfill

\begin{center}
    An electronic version of this thesis is available at http://repository.tudelft.nl/.
\end{center}

\twocolumn
\let\clearpage\clearpagebackup  
\clearpage
\setcounter{page}{1}

\onecolumn

\begin{abstract}
The aim of this template is to make it more clear what is expected from you. 
\textbf{It is by no means required to follow this exact same structure.}
The abstract should be short and give the overall idea:
what is the background, the research questions, what are your contributions, and what are the main conclusions.
It should be readable as a stand-alone text (preferably no references to the paper or to outside literature).
\end{abstract}

\section{Introduction}
The transformer architecture invention led to the rise of various language models. One of the challenges the researchers have encountered with such models is the speed of autoregressive inference. Optimizing the inference is important whether aiming to provide a good user experience in real-time applications, saving on computational costs when running the model, or allowing running the model in an environment with limited computational and memory resources. 

In particular, memory bandwidth is a significant limitation in modern highly parallelizable hardware \cite{pope_efficiently_2022}\cite{shazeer_fast_2019}\cite{williams_roofline_2009}. In incremental decoding, most of the bandwidth is consumed by the constant loading of large key and value tensors \cite{shazeer_fast_2019}. To address this issue, Multi-Query Attention (MQA) architecture was proposed and was shown to increase the inference speed significantly though the generation quality was diminished \cite{shazeer_fast_2019}. To address the quality decline that MQA induces, Grouped-Query Attention (GQA) was introduced as an interpolation between the standard Multi-Head Attention (MHA) and MQA \cite{ainslie_gqa_2023}.

While Ainslie et al. \cite{ainslie_gqa_2023} present the GQA models' performance and speed results for different query groups, the models used in the experiments were pre-existing and then further uptrained. This research seeks to search for an optimal number of query groups relative to the number of attention heads at the pre-training stage to utilize the full potential of GQA. Moreover, the experiments conducted in this study focus on increasingly popular decoder-only models, where such optimizations play the most noticeable role due to the models' autoregressive nature. 

Besides, we are testing an orthogonal approach to reducing the size of key and value tensors, namely - by shrinking the dimensionality of key/query, and value vectors (kqv-vectors) while keeping the number of attention heads the same, to see the direct effect of their dimensionality on speed and quality. Even though Shazeer \cite{shazeer_fast_2019} has introduced this idea, the experiments have not revealed the direct correlation between the size of kqv-vectors and the models' quality, as with their size the number of attention heads was adjusted respectively. In addition, in those experiments for kqv-vectors' size, the speed measurements were omitted fully. 

Hence, in this study, we first try to find an optimal dimensionality for key/query and value vectors while keeping the number of attention heads the same. Next, we run GQA experiments with different proportions of groups relative to the total number of attention heads. In addition, once the optimal proportion of groups is determined, we investigate whether the quality and speed ratio relative to the baseline is preserved when adjusting the number of attention heads while keeping the proportion of groups constant. Finally, we are testing the combination of both approaches by setting the most optimal kqv-vectors' size in the GQA model with the best trade-offs. 

\begin{itemize}
    \item Main findings and how they contribute, and mention again that the goal is to find optimal trade-offs between speed and quality.
\end{itemize}

% \begin{itemize}
% \item Relate to the most relevant existing work from the literature (use BibTeX), explain their contributions, and (critically) indicate what is still unanswered. 
% %\emph{The existing state of the art describes the setup of general scientific papers, e.g.\ see~\cite{hengl2002rules}, but this may be different for computer science papers.}

% \item Explain what the research questions for this work are. 
% This usually is a subset of the unanswered questions. %\emph{The aim of this work is therefore to provide a good template specifically for papers in the field of computer science.}

% \item Summarize the main contributions/conclusions of this research.
% NB: Make sure the title of the paper is a good match to the main research question / contribution / conclusion.

% \item Briefly indicate how the rest of the paper fits together to answer the research question(s).
% \end{itemize}

% For a longer research paper, a section with a more elaborate discussion of the literature may follow, but for short (conference) submissions, this is often included in the introduction.

% Make sure the introduction and conclusion are easily understandable by everyone with a computer science bachelor (e.g.\ your examiner may have a completely different expertise).

\section{Literature Review and Problem Statement}
\subsection{Background}
Transformers is a state-of-the-art deep learning architecture in large language models used for data processing (encoding) and data generation tasks (decoding) \cite{vaswani_attention_2017}. Over time, transformer models have been scaled to enormous sizes with hundreds of billions of parameters to improve performance. This made the models less sustainable and more expensive both to train and operate due to the computational requirements induced by the scaling. 

For this reason, the TinyStories dataset has specifically been designed to demonstrate the capability of improving the models' sample efficiency. It has been shown that improving the sample efficiency can significantly reduce the number of parameters while preserving the model's ability to generate coherent, diverse, and consistent results with the model demonstrating some reasoning capabilities \cite{eldan_tinystories_2023}. To facilitate the sample efficiency in models, the TinyStories was designed to be more cohesive and suitable for learning English grammar and vocabulary.

Furthermore, the research reveals that the smaller models exhibit similar behaviour to their larger counterparts when it comes to architectural decisions \cite{eldan_tinystories_2023}. This, coupled with their suitability for limited-resources environments, positions them as ideal platforms for architectural experimentation in scenarios with restricted access to data and computational resources. 

\subsection{Problem Definition}
\textit{Introduce the importance of inference speed and what's its current bottleneck.}

\subsection{Related Work}
\textit{Introduce MQA, GQA, and gaps for both.}

% For some types of work in computer science the methodology is standard: analyze the problem (e.g., make assumptions and derive properties), present a new algorithm and its theoretical background, proving its correctness, and evaluate unproven aspects in simulation.
% Then an explanation of the methodology is often omitted, and the setup of the evaluation is part of a later section on the evaluation of the ideas.\footnote{This already shows that there is no single outline to be given for all papers.}
% In this case, explain relevant (background) concepts, theory and models in this section (with references) and relate them to your research question.
% Also this section then typically contains a more precise, formal description of the problem.

% Do not forget to give this section another name, for example after the problem you are solving.


\section{Approach}
% In computer science typically the third section contains an exposition of the main ideas, for example the development of a theory, the analysis of the problem (some proofs), a new algorithm, and potentially some theoretical analysis of the properties of the algorithm.

% Do not forget to give this section another name, for example after the method or idea you are presenting.

% Some more detailed suggestions for typical types of contributions in computer science are described in the following subsections.

This research will test different architectural decisions for the attention mechanism in transformer models via the TinyStories dataset. To achieve this, small models such as RoBERTa\footnotemark[1]{} and GPT-Neo\footnotemark[2]{} will be used to test the decisions while covering both text processing and text generation types of tasks. Then, the performance of both models for different architectural decisions will be evaluated with different metrics included in the BabyLM evaluation pipeline\footnotemark[3]{}.

The architectural decisions tested in this research will be the multi-query attention mechanism, the grouped-query mechanism, and the standard multi-head mechanism with a reduced size of query/key vectors. All of them are expected to increase the inference speed due to a reduction in memory overhead \cite{ainslie_gqa_2023} \cite{shazeer_fast_2019}. 

\textit{Introduce the time complexity analysis and how the memory affects it. Also, mention that we are testing at the pre-training stage. And, mention why it could be beneficial to have it done at the pre-training stage. Maybe also mention the changing number of attention heads as a gap from the previous research.}


\subsection*{Experimental work}
In this case, this section will mostly contain a description of the methods/algorithms you will be comparing. Although not all methods need to be described in detail (providing appropriate references are available), make sure that you reveal sufficient details to a reader not familiar with these methods to: a) obtain a high-level understanding of the method and differences between them, and b) understand your explanation of the results/conclusions.

\subsection*{Improvement of an idea}
In this case, you would need to explain in detail how the improvement works. If it is based on some observation that can be proven, this is a good place to provide that proof (e.g., of the correctness of your approach). 

\subsection*{Literature survey}
If your contribution is a literature survey, then the organization of these ``middle'' sections very much depends on the way you want to present/organize the literature you are discussing.
First try to cluster papers that are similar in some aspect. Then think of how these clusters are related, from that you can think of a good order to discuss these clusters; this is sometimes called a bottom-up approach to writing a paper.

In addition, you may try to think about the organization of the literature from a top-down perspective: try to ``take a step back'' and think about the field and what important questions/variants are and build a hierarchical categorization of the field.

Make clear what your contribution is here: a new organization of the literature, identification of open problems/challenges, new parallels/generalizations, a table with pros/cons of different methods, etc.\ 



%\bigskip
%\lipsum[1-67]

\section{Experiments and Results}
\subsection{Experiments and Their Order}
\begin{itemize}
    \item Explicitly introduce the research questions
    \item Explain how the experiments were conducted
    \item Add a nice diagram for the order
\end{itemize}
\subsection{Experimental Setup}
\subsection{Models' Quality}
\subsection{Models' Speed}


\section{Responsible Research}
Reflect on the ethical aspects of your research and discuss the reproducibility of your methods.
Note that although in many published works there is no such a section (it may be part of some meta-information collected by the journal, or part of the discussion section), we require you to think (and report) about this as part of this course.

\section{Discussion}
Results can be compared to known results and placed in a broader context.
Provide a reflection on what has been concluded and how this was done.
Then give a further possible explanation of results.


\section{Conclusions and Future Work}
Briefly summarize the (main) research question(s).
Provide your conclusions, the answers to the research question(s).
Make statements.
Highlight interesting elements, contributions.

Discuss open issues, possible improvements, and new questions that arise from this work; formulate recommendations for further research.

Ideally, this section can stand on its own: it should be readable without having read the earlier sections and accessible to anyone with a bachelor degree in Computer Science.


\appendix
\section{Some further guidelines that go without saying (right?)}

\begin{itemize}
\item Read the manual for the Research Project. (See e.g.\ the instructions on the maximum length: less is more!)
\end{itemize}

\subsection{Reference use}
\begin{itemize}
\item use a system for generating the bibliographic information automatically from your database, e.g., use BibTex and/or Mendeley, EndNote, Papers, or \ldots
\item all ideas, fragments, figures and data that have been quoted from other work have correct references
\item literal quotations have quotation marks and page numbers
\item paraphrases are not too close to the original
\item the references and bibliography meet the requirements
\item every reference in the text corresponds to an item in the bibliography and vice versa
\end{itemize}

\subsection{Structure}
Paragraphs
\begin{itemize}
\item are well-constructed
\item are not too long: each paragraph discusses one topic
\item start with clear topic sentences
\item are divided into a clear paragraph structure
\item there is a clear line of argumentation from research question to conclusions
\item scientific literature is reviewed critically
\end{itemize}

\subsection{Style}
\begin{itemize}
\item correct use of English: understandable, no spelling errors, acceptable grammar, no lexical mistakes 
\item the style used is objective
\item clarity: sentences are not too complicated (not too long), there is no ambiguity
\item attractiveness: sentence length is varied, active voice and passive voice are mixed
\end{itemize}

\subsection{Tables and figures}
\begin{itemize}
\item all have a number and a caption
\item all are referred to at least once in the text
\item if copied, they contain a reference
\item can be interpreted on their own (e.g. by means of a legend)
\end{itemize}


\bibliographystyle{plain}
\bibliography{references}

A rule of thumb for dealing with the literature is the following: scan about 10--20 contributions: read title, abstract, part of introduction and conclusions; categorize contribution; some of these are studied in more depth: completely read about 5 conference papers or equivalent (summarize contribution in own words); of which studied in-depth about 2 conference papers (the student is able to explain in detail and criticize contributions). This may result in 5--20 references, possibly even more if the project is a literature study.

\end{document}
